\documentclass[10pt,a4paper]{amsart}
\usepackage{amsmath,amssymb,amsthm,graphicx}
\usepackage{verbatim}
\newenvironment{reference}{\comment}{\endcomment}
%\newenvironment{reference}{}{}
\newenvironment{slogan}{\comment}{\endcomment}
\newenvironment{history}{\comment}{\endcomment}

% For commutative diagrams we use Xy-pic
\usepackage[all]{xy}

% We use 2cell for 2-commutative diagrams.
\xyoption{2cell}
% Theorem environments.
%
\theoremstyle{plain}
\newtheorem{theorem}[subsection]{Theorem}
\newtheorem{proposition}[subsection]{Proposition}
\newtheorem{lemma}[subsection]{Lemma}

\theoremstyle{definition}
\newtheorem{definition}[subsection]{Definition}
\newtheorem{example}[subsection]{Example}
\newtheorem{exercise}[subsection]{Exercise}
\newtheorem{situation}[subsection]{Situation}

\theoremstyle{remark}
\newtheorem{remark}[subsection]{Remark}
\newtheorem{remarks}[subsection]{Remarks}

\numberwithin{equation}{subsection}

% Macros
%
\def\lim{\mathop{\mathrm{lim}}\nolimits}
\def\colim{\mathop{\mathrm{colim}}\nolimits}
\def\Spec{\mathop{\mathrm{Spec}}}
\def\Hom{\mathop{\mathrm{Hom}}\nolimits}
\def\Ext{\mathop{\mathrm{Ext}}\nolimits}
\def\SheafHom{\mathop{\mathcal{H}\!\mathit{om}}\nolimits}
\def\SheafExt{\mathop{\mathcal{E}\!\mathit{xt}}\nolimits}
\def\Sch{\mathit{Sch}}
\def\Mor{\mathop{\mathrm{Mor}}\nolimits}
\def\Ob{\mathop{\mathrm{Ob}}\nolimits}
\def\Sh{\mathop{\mathit{Sh}}\nolimits}
\def\NL{\mathop{N\!L}\nolimits}
\def\CH{\mathop{\mathrm{CH}}\nolimits}
\def\proetale{{pro\text{-}\acute{e}tale}}
\def\etale{{\acute{e}tale}}
\def\QCoh{\mathit{QCoh}}
\def\Ker{\mathop{\mathrm{Ker}}}
\def\Im{\mathop{\mathrm{Im}}}
\def\Coker{\mathop{\mathrm{Coker}}}
\def\Coim{\mathop{\mathrm{Coim}}}

% Boxtimes
%
\DeclareMathSymbol{\boxtimes}{\mathbin}{AMSa}{"02}

%
% Macros for moduli stacks/spaces
%
\def\QCohstack{\mathcal{QC}\!\mathit{oh}}
\def\Cohstack{\mathcal{C}\!\mathit{oh}}
\def\Spacesstack{\mathcal{S}\!\mathit{paces}}
\def\Quotfunctor{\mathrm{Quot}}
\def\Hilbfunctor{\mathrm{Hilb}}
\def\Curvesstack{\mathcal{C}\!\mathit{urves}}
\def\Polarizedstack{\mathcal{P}\!\mathit{olarized}}
\def\Complexesstack{\mathcal{C}\!\mathit{omplexes}}
% \Pic is the operator that assigns to X its picard group, usage \Pic(X)
% \Picardstack_{X/B} denotes the Picard stack of X over B
% \Picardfunctor_{X/B} denotes the Picard functor of X over B
\def\Pic{\mathop{\mathrm{Pic}}\nolimits}
\def\Picardstack{\mathcal{P}\!\mathit{ic}}
\def\Picardfunctor{\mathrm{Pic}}
\def\Deformationcategory{\mathcal{D}\!\mathit{ef}}
\usepackage{cleveref}
\crefname{originallemma}{original lemma}{original lemmas}
\Crefname{originallemma}{Original lemma}{Original lemmas}
\makeatletter
\newlabel{lemma-duality-proper-over-field}{{0FVV}{source}{}{sourcepointer-1}{}}
\newlabel{lemma-duality-proper-over-field@cref}{{[originallemma][0][]0FVV}{[1][0][]source}{}{}{}}
\makeatother
\crefname{originalreference}{original reference}{original references}
\Crefname{originalreference}{Original reference}{Original references}
\makeatletter
\newlabel{remark-duality-proper-over-field}{{0FVW}{source}{}{sourcepointer-2}{}}
\newlabel{remark-duality-proper-over-field@cref}{{[originalreference][0][]0FVW}{[1][0][]source}{}{}{}}
\makeatother
\crefname{originalsection}{original section}{original sections}
\Crefname{originalsection}{Original section}{Original sections}
\makeatletter
\newlabel{cohomology-section-cup-product}{{0FKU}{source}{}{sourcepointer-3}{}}
\newlabel{cohomology-section-cup-product@cref}{{[originalsection][0][]0FKU}{[1][0][]source}{}{}{}}
\makeatother
\crefname{originalsection}{original section}{original sections}
\Crefname{originalsection}{Original section}{Original sections}
\makeatletter
\newlabel{cohomology-example-dual-derived}{{0FPC}{source}{}{sourcepointer-4}{}}
\newlabel{cohomology-example-dual-derived@cref}{{[originalsection][0][]0FPC}{[1][0][]source}{}{}{}}
\makeatother
\crefname{originallemma}{original lemma}{original lemmas}
\Crefname{originallemma}{Original lemma}{Original lemmas}
\makeatletter
\newlabel{cohomology-lemma-ext-composition-is-cup}{{0FVB}{source}{}{sourcepointer-5}{}}
\newlabel{cohomology-lemma-ext-composition-is-cup@cref}{{[originallemma][0][]0FVB}{[1][0][]source}{}{}{}}
\makeatother
\title{Perfect pairings for a perfect complex against the dualizing complex of a proper scheme over a field}
\author{The Stacks Project Authors}
\date{stacks-project:duality.tex:local-snapshot-20261012; GNU FDL 1.2 or later}
\begin{document}
\maketitle
\noindent\textit{Editorial note.} All delivered mathematical text is the author's own, copied from the pinned Stacks Project chapter file duality.tex; no statement or proof step was rewritten, reordered or removed.\par\medskip
\noindent\textit{Editorial note.} Preamble provenance: the pinned chapter file opens with an input of the shared preamble and carries no preamble text of its own, so the project's single Stacks-wide preamble staged at sources/stacks-duality/preamble.tex is excerpted verbatim under the recipe fixed by LANE\_BRIEF 22.1, as three spans, lines 8-13 (verbatim package and the reference/slogan/history environments), lines 14-18 (Xy-pic setup) and lines 154-217 (the native theorem declarations and mathematical macros). Nothing was added, removed or redefined; the wrapper class and the amsmath/amssymb/amsthm/graphicx base packages are supplied by the assembler.\par\medskip
\noindent\textit{Editorial note.} The counted statement and proof contain no cross reference, no citation and no included file, so the delivered body is self contained.\par\medskip
\noindent\textit{Editorial note.} Canonical location of the counted proposition: Stacks tag 0FVY (sources/stacks-duality/pinned-original-source.tex line 7078).\par\medskip
\noindent\textit{Allowed standard prerequisites (editorial).}\begin{itemize}
\item Duality for a proper scheme over a field. Same chapter, Lemma labelled lemma-duality-proper-over-field, cited at proof lines 7097-7099. Fixes k, X and the dualizing complex.
\item The trace pairing on a proper scheme over a field. Same chapter, Remark labelled remark-duality-proper-over-field, cited at proof lines 7097-7103. Contains the pairing whose perfectness is invoked.
\item Composition of Ext groups is given by the cup product. Cohomology, Lemma labelled cohomology-lemma-ext-composition-is-cup, cited at proof lines 7099-7101. Identifies the displayed pairing with the one of the remark.
\end{itemize}\medskip
\ifdefined\hypertarget\hypertarget{sourcepointer-1}{}\fi
\noindent\textit{Source reference.} Original lemma 0FVV: sources/stacks-duality/pinned-original-source.tex, source lines 6934--6934. Same-chapter pointer into the author environment lemma at source line 6934; invoked at proof line None of the counted proof. It is retained only as a navigation pointer; every load-bearing step is inside the delivered span.\par\medskip
\ifdefined\hypertarget\hypertarget{sourcepointer-2}{}\fi
\noindent\textit{Source reference.} Original reference 0FVW: sources/stacks-duality/pinned-original-source.tex, source lines 7011--7011. Same-chapter pointer into the author environment remark at source line 7011; invoked at proof line 7101 of the counted proof. It is retained only as a navigation pointer; every load-bearing step is inside the delivered span.\par\medskip
\ifdefined\hypertarget\hypertarget{sourcepointer-3}{}\fi
\noindent\textit{Source reference.} Original section 0FKU: sources/stacks-chow/tags.txt, source lines 19474--19474. Cross-chapter standard prerequisite located in the project tag list staged in the local snapshot; invoked at proof line None of the counted proof. It is retained only as a navigation pointer; every load-bearing step is inside the delivered span.\par\medskip
\ifdefined\hypertarget\hypertarget{sourcepointer-4}{}\fi
\noindent\textit{Source reference.} Original section 0FPC: sources/stacks-chow/tags.txt, source lines 19597--19597. Cross-chapter standard prerequisite located in the project tag list staged in the local snapshot; invoked at proof line None of the counted proof. It is retained only as a navigation pointer; every load-bearing step is inside the delivered span.\par\medskip
\ifdefined\hypertarget\hypertarget{sourcepointer-5}{}\fi
\noindent\textit{Source reference.} Original lemma 0FVB: sources/stacks-chow/tags.txt, source lines 19806--19806. Cross-chapter standard prerequisite located in the project tag list staged in the local snapshot; invoked at proof line 7099 of the counted proof. It is retained only as a navigation pointer; every load-bearing step is inside the delivered span.\par\medskip
\begin{lemma}
\label{lemma-duality-proper-over-field-perfect}
Let $k$, $X$, and $\omega_X^\bullet$
be as in Lemma \ref{lemma-duality-proper-over-field}.
Let $t : H^0(X, \omega_X^\bullet) \to k$ be as in
Remark \ref{remark-duality-proper-over-field}.
Let $E \in D(\mathcal{O}_X)$ be perfect. Then the pairings
$$
H^i(X, \omega_X^\bullet \otimes_{\mathcal{O}_X}^\mathbf{L} E^\vee)
\times
H^{-i}(X, E) \longrightarrow k, \quad
(\xi, \eta) \longmapsto
t((1_{\omega_X^\bullet} \otimes \epsilon)(\xi \cup \eta))
$$
are perfect for all $i$. Here $\cup$ denotes the cup product
of Cohomology, Section \ref{cohomology-section-cup-product} and
$\epsilon : E^\vee \otimes_{\mathcal{O}_X}^\mathbf{L} E \to \mathcal{O}_X$
is as in Cohomology, Example \ref{cohomology-example-dual-derived}.
\end{lemma}

\begin{proof}
By replacing $E$ with $E[-i]$ this reduces to the case $i = 0$.
By Cohomology, Lemma \ref{cohomology-lemma-ext-composition-is-cup}
we see that the pairing is the same as the one discussed
in Remark \ref{remark-duality-proper-over-field}
whence the result by the discussion in that remark.
\end{proof}

\end{document}
